\documentclass[11pt]{article}
\usepackage{amsmath,amssymb,hyperref}
\title{Fallible probes: physical cleanup, timely evidence and small rejection certificates}
\author{}\date{}
\begin{document}\maketitle
Two civil inspection crew roles have separate observations. The
scout knows the relevant local condition $x$ and the receiving
crew knows its own route condition $y$. The tight repair uses
route $x\oplus y$. A usable observation conveys the actual
upstream condition, before the receiver commits. If two
histories with the same local $y$ produce no usable read and
different $x$, that crew has no one compatible repair.
A more accurate read that is private or late leaves the same
obstruction. Physical mechanisms below supply or remove this
evidence; it is not an oracle at the receiver.

\paragraph{Charge physical attempts and old work.}
At most $k$ acquisition attempts can be erased on an admitted trace.
An ordinary shared-port attempt takes one tick. A shared-port
failure creates a $d$-tick physical cleanup, and the old duty
needs a physically staffed standby throughout. Without the standby,
an attempt that might fail is inadmissible. An isolated-port attempt
takes $s\ge1$ ticks and spends $c>0$ battery units at launch, even
when it returns no observation. Its failure does not occupy the
shared cleanup resource. Both modes are assumed to report a
true value on successful acquisition, under the declared source
law, not measured to do so universally.

For a serial facility, all drive and cleanup phases are sequential.
Given an already staffed standby, if $j$ of the $k$ failed attempts
are isolated, the last guaranteed-success attempt can be shared for
one tick. The worst completion is
$1+j s+(k-j)(1+d)$, with $0\le j\le\min(k,\lfloor B/c\rfloor)$.
The minimum uses as many isolated trials as allowed iff $s<1+d$.
Without a standby, all $k$ potentially failing trials must be isolated.
This is necessity as well as sufficiency: before a first real read
the policy has one no-cue history, and the admitted law can extend it
for $k$ failures. A policy may stop earlier on success, but cannot
count an unobserved proposal or a failed battery expenditure as a read.

If a separately staffed isolated port can operate while physical
cleanup continues, the old duty's deadline and the new observation's
deadline are different. Track current time, remaining cleanup,
battery and erasures. Starting a shared attempt waits for the actual
end of cleanup, while an isolated one can start immediately. No
attempt overlaps the previous acquisition; it can overlap other work.
Early success does not settle the old cleanup. With $d=3,s=2,k=2$
and enough battery for one isolated trial, the worst schedule is:
failed shared $[0,1)$; old cleanup $[1,4)$; isolated $[1,3)$;
then shared $[4,5)$. Thus the exact evidence deadline falls
from seven to five while old cleanup must still end by four.
If the isolated test succeeds at three, old cleanup is still live.
The serialized and concurrent source roles differ in physical staffing,
not a matched algorithm advantage.

A closed common-horizon test also follows for a homogeneous source.
For $r$ failed shared attempts, each cleanup block can contain
$a=\lfloor d/s\rfloor$ isolated trials without extension.
Its next trial costs $s-b$ extra time, $b=d\bmod s$, and subsequent
trials cost $s$. If $J$ isolated trials are assigned after any of
the $r$ failed shared attempts, the minimum extra time is
$(s-b)\min(q,r)+s\max(q-r,0)$ for $q=\max(J-ra,0)$; with no shared
attempts this gives $sJ$. Minimize over the split of the $k$ failures
and final probe type. The base cost is $r(1+d)$ and a final shared
attempt costs one additional tick. For a final isolated probe,
count it in $J$ and place it in the last cleanup block if any.
Before the first shared failure, isolated work can instead be
moved into some cleanup block without raising this common horizon.
Filling the cheapest per-block increments proves optimality.
This formula is for a common deadline; the separate old and new
deadlines use the event solver rather than collapsing to their max.

\paragraph{One fixed failure law, not independent maxima.}
If a latent type specifies two budgets, one per physical port,
decrement only the failed port's budget in each surviving type and
remove types that could not erase it. For this failure-only source
support, discard dominated budget vectors: a dominated type can
generate no additional failure trace. Either decrement preserves
the pointwise order, so the maximal frontier remains exact at
each no-cue step. This does not permit forgetting other physical
or decision facts that differ between two types.

For coupled budgets $(1,0)$ or $(0,1)$, without an extra standby,
a two-tick isolated try can fail, but genuine isolated failure then
implies the shared try cannot fail. One isolated charge yields an
exact worst deadline of three. Replacing that source support with
the independent maximum $(1,1)$ breaks this guarantee at the same
battery budget. Absence of remote delivery is not automatically the
primitive failed-acquisition event used in that proof. One trace
can have no local record and safe shared retry; another may retain
a true private record but have a lost transmission, while shared
retry can still fail and would disrupt the old duty. With those
source histories and one ambiguous no-cue observation, neither
follow-up is universally correct. A separately charged local record
check followed by an authorized physical retransmission can restore
service if the cutoff leaves room and the replay resource is actually
free. A record seen by a sender alone does not furnish a receiver's
route before delivery.

\paragraph{A cutoff witness for live assignment.}
For one receiving observation class, put an edge from each of $n$
duties to each of $m$ single-use slots only if the assignment is
lawful and physically executable in all its source histories.
For $1\le n\le m$ a missing complete matching has a rejecting
subfamily of size at most
\[
             \max_{1\le t\le n} t(m-t+1).
\]
A Hall cut with $t$ duties leaves at least $m-t+1$ slots outside
its common neighborhood. Take that many and one excluding
history for each of their pairs with those duties. Those
histories still reject the assignment. For sharpness, choose
$t$ duties and $t-1$ common slots. Add one history deleting
exactly each specific pair of a chosen duty and any other
slot and no other edges. All histories are necessary: removing
one returns an edge, connects one selected duty, and leaves
enough slots for all other duties. This is a certificate for a
fixed observation class and an authorized simultaneous
assignment, not a live sensor that can discover its own
source histories. The proof is closely related to the established
\href{https://doi.org/10.1002/jgt.20180}{characterization of matching blockers};
the distinction here is grouping edge exclusions by physical
source history. No historical novelty of the combinatorial
bound is asserted.

The physical two-duty control derives a coarse read from
continuity across each pair of physical contacts. The
first partial reading survives timeout at the same
declared receiving cut. A timeout on the shared feed also
holds slot S for cleanup, exhausting remaining common slots.
An otherwise matched independently powered trial preserves S
and makes that partial read sufficient. Private or late exact
fault information does not help. The instantiated contacts,
power and timing are modeled source primitives, not measured
hardware.

\paragraph{Comparison and evidence.}
Primary baselines: \href{https://arxiv.org/abs/1209.1695}{common-information dynamic programming},
\href{https://arxiv.org/abs/1410.0083}{reactive active sensing},
\href{https://drops.dagstuhl.de/storage/00lipics/lipics-vol391-concur2026/html/LIPIcs.CONCUR.2026.37/LIPIcs.CONCUR.2026.37.html}{active diagnosis with costs and rewards},
and \href{https://doi.org/10.1112/jlms/s1-10.37.26}{Hall's matching theorem}.
Charging sensing or using dynamic beliefs is established.
Same-source direct policy enumeration and matching tie in every
tested scope. Replays include 8,640 serial, 12,000 overlap,
4,480 common-horizon, 3,500 coupled-law, 120 receipt and
9,000 matched Hall controls with 165 sharp constructions.
Four terms check in Lean. This is a constructive and obstructive
R3a reference with a declared model; it does not identify the true
latent source law from data (R3b), independently validate a mission,
or establish advancement against the strongest fully matched planner.
\end{document}
